\documentclass[11pt]{article}
\usepackage{amsmath, amssymb, amsfonts}
\usepackage{geometry}
\geometry{margin=2.5cm}

\title{Perfektoide Zahlen als lineare Optimierungsprobleme}
\author{}
\date{}

\begin{document}
\maketitle

\section*{Vorabklärung}

Es gibt \emph{keine kanonische mathematische Äquivalenz} zwischen perfektoiden Zahlen und linearen Optimierungsproblemen.
Was hier beschrieben wird, ist eine \emph{strukturtreue Modellierung}:
Eine perfektoide bzw.\ p-adische Zahl wird in eine endliche Menge linearer Variablen,
eine Hauptbedingung und mehrere Nebenbedingungen übersetzt,
sodass ein Simplex- oder MILP-Verfahren darauf anwendbar ist.

\section*{1. p-adische und perfektoide Struktur}

Eine p-adische Zahl besitzt eine Darstellung
\[
x = \sum_{k \ge k_0} a_k p^k,
\qquad a_k \in \{0,1,\dots,p-1\}.
\]

Die zugehörige p-adische Bewertung ist
\[
v_p(x) = \min\{k \mid a_k \neq 0\}.
\]

Perfektoide Strukturen sind zusätzlich durch die Surjektivität des Frobenius
\[
\varphi(x) = x^p
\]
charakterisiert.

Die Geometrie ist nicht-archimedisch:
\[
|x+y|_p \le \max(|x|_p, |y|_p).
\]

\section*{2. Variablenwahl für das Optimierungsproblem}

Wir fixieren eine endliche Approximation
\[
x \approx \sum_{k=k_0}^{k_1} a_k p^k.
\]

Daraus entstehen die Optimierungsvariablen
\[
\mathbf{a} = (a_{k_0}, a_{k_0+1}, \dots, a_{k_1}) \in \mathbb{R}^{n}.
\]

\section*{3. Hauptbedingung (Zielfunktion oder Gleichung)}

\subsection*{Variante A: Rekonstruktion einer Zielzahl}

Für einen gegebenen Zielwert \(X\):
\[
\sum_{k=k_0}^{k_1} a_k p^k = X.
\]

Dies ist eine lineare Gleichung und kann als Hauptbedingung dienen.

\subsection*{Variante B: Minimierung der p-adischen Bewertung}

Die p-adische Tiefe wird linear approximiert durch
\[
\min \sum_{k=k_0}^{k_1} k \cdot a_k.
\]

Diese Zielfunktion bevorzugt niedrige Bewertungen
und modelliert p-adische ``Kleinheit''.

\section*{4. Nebenbedingungen}

\subsection*{(1) Ziffernbeschränkung}

Die p-adische Ziffernstruktur erzwingt
\[
0 \le a_k \le p-1
\quad \text{für alle } k.
\]

\subsection*{(2) Frobenius-Kohärenz (perfektoide Approximation)}

Perfektoidheit verlangt die Existenz eines p-ten Wurzelelements.
Diskret modelliert:
\[
a_{k+p} \approx a_k.
\]

Linear relaxiert:
\[
a_{k+p} - a_k = 0
\quad \text{oder} \quad
|a_{k+p} - a_k| \le \varepsilon.
\]

\subsection*{(3) Ultrametrische Dominanz}

Für eine Summe \(x = y + z\) gilt p-adisch:
\[
v_p(x) = \min(v_p(y), v_p(z)).
\]

Linear relaxiert ergibt sich:
\[
a_k \le b_k + c_k.
\]

\subsection*{(4) Endlichkeits- und Trunkierungsbedingung}

Zur Sicherstellung eines endlichen Problems:
\[
\sum_{k=k_0}^{k_1} a_k \le M,
\]
wobei \(M\) eine feste Schranke ist.

\section*{5. Vollständiges lineares Optimierungsproblem}

\paragraph{Variablen:}
\[
a_{k_0}, \dots, a_{k_1}.
\]

\paragraph{Zielfunktion (Beispiel):}
\[
\min \sum_k k a_k.
\]

\paragraph{Nebenbedingungen:}
\begin{align*}
0 &\le a_k \le p-1, \\
a_{k+p} - a_k &= 0, \\
\sum_k a_k &\le M.
\end{align*}

Dieses Problem ist ein lineares Programm (LP)
oder bei Ganzzahligkeit der \(a_k\) ein MILP,
und damit direkt vom Simplex-Algorithmus behandelbar.

\section*{6. Interpretation}

\begin{itemize}
\item Es wird nicht über perfekte Zahlen optimiert,
      sondern über ihre endlichen Darstellungen.
\item Die p-adische und perfektoide Struktur
      erscheint als Nebenbedingungsgeometrie.
\item Die Abbildung ist keine Isomorphie,
      sondern eine strukturtreue Relaxierung.
\end{itemize}

\section*{7. Fazit}

Perfektoide Zahlen können sinnvoll als
lineare Optimierungsprobleme modelliert werden,
wenn man ihre Bewertung, Ziffernstruktur und Frobenius-Eigenschaften
in Zielfunktionen und Nebenbedingungen übersetzt.
Der Simplex-Algorithmus optimiert dabei nicht Zahlen,
sondern ihre algebraisch motivierten Repräsentationen.

\end{document}

